\chapter*{Abstract}
\bigskip


How does brain connectivity shape computation? Although connectome-constrained recurrent networks are widely used to link structural organisation to function, previous studies have largely focused on network performance, leaving unclear how connectivity shapes the geometry of neural activity and how this geometry determines computational capacity. Here we show that the connectome's key advantage is not greater peak capacity, but a wider range of activity levels over which computation remains possible. We compare a human structural connectome with matched null networks, quantifying computational capacity through two complementary tasks: memory retention and autonomous signal prediction. We find that where its strong connections sit keeps activity from collapsing into a single population-wide pattern. The connectome keeps computing as activity is pushed well past the level at which the nulls have already collapsed. What survives that collapse is a set of directions varying independently of one another, and it is these that support memory. Conventional variance-based dimensionality, such as participation ratio, discounts these directions by construction, whereas an effective rank counting directions by their contribution to what can be read out closely tracks memory capacity. Prediction imposes an opposing constraint: alignment with this population-wide pattern simplifies trajectories and improves prediction, but reduces the independent dimensions required for memory. Thus, connectivity shapes computation by controlling the geometry of neural dynamics rather than by maximising capacity, identifying activity-manifold geometry as a mechanistic link between structure and function.  

\bigskip
\bigskip


\noindent\textbf{Keywords:} brain-computer interface (BCI), electroencephalography (EEG), motor imagery (MI), classification, machine learning, deep learning, human-robot interaction, neuroengineering.